\section{Resource Outcomes in the Economy}
\label{sec:outcomes}

\subsection{Electricity cost}

When a resource plan raises the cost of electricity, two groups pay. In the Philippines,
industries buy nearly three quarters of the electricity used in the economy, so most of
the cost falls on production. It raises firms' costs directly and through the inputs they
buy from one another, and it reaches households through prices and wages. Households pay
most of the remainder in their own bills. Because every household needs a minimum amount
of electricity, the same increase takes a larger share of a poorer household's budget.
OG--CLEWS carries the cost to both groups. The household share is represented as a higher
price of electricity, designed so that it does not create revenue for government: a higher
cost of electricity is a cost to the economy, not a tax.

In the Philippine example, the scenario raises the cost of electricity by about 5 percent
on average, and by about 16 percent at its peak around 2030. Figure~\ref{fig:electricity}
shows where that cost lands. Applied to household bills alone, even the full increase
barely moves GDP, although consumption is about 0.06 percent lower over the first decade.
When the cost also reaches industry, GDP is about 0.19 percent lower over the first
decade, nearly half a percent lower at the 2030 peak, and still 0.14 percent lower in 2075.
Most of the economy-wide effect therefore comes through production, not household bills.

\begin{figure}[H]
\centering
\includegraphics[width=\linewidth]{electricity-routes.pdf}
\caption{Electricity cost: GDP and consumption when the price increase reaches household
bills only, and when it also raises industry costs. Percent change from the control,
Philippine example. The household-only case applies the full increase to the combined
energy-and-water good without returning any revenue, a stronger household shock than in
the combined case.}
\label{fig:electricity}
\end{figure}

\subsection{Public infrastructure}

A resource plan may require new public infrastructure, such as grid and transmission
lines. Where these assets are publicly financed, their cost enters government investment
in OG-Core. Public capital raises the productivity of the economy, but it must be paid for
through taxes or debt, and debt financing can displace private investment. OG--CLEWS
follows both sides: the gain in productivity and the cost of financing it. Which assets are
public, and how they are financed, is specified for each scenario rather than inferred from
the type of asset.

In the Philippine example, the scenario mainly shifts public spending between years rather
than adding a large programme. Annual public investment moves between about 1.3 percent
below and 0.4 percent above the control, and the largest change in GDP is below 0.002
percent. A scenario with a substantial publicly financed grid expansion would engage this
channel more fully.

\subsection{Generation capital intensity}

A cleaner generation mix is usually more capital-intensive. Solar and wind plants are paid
for almost entirely when they are built, while fossil-fuel plants cost less to build but
buy fuel throughout their lives. OG--CLEWS carries this change into the economy by raising
the share of capital in the production of electricity. The result is not the one intuition
suggests. In the Philippine calibration, a larger capital share lowers the cost of
producing electricity, so its price falls and output rises. But the industry's revenue falls
with the price, and an industry holds capital in proportion to its revenue. The size of
these responses depends on the calibration; the distinction between technology and
investment does not.

In the Philippine example, capital's share in electricity production rises from 78 to 80
percent. In the long run, the price of electricity falls by 11 percent and output rises by
4.4 percent, but the industry's revenue falls by 7 percent. Its capital falls by 4.8
percent and its labour by 18.7 percent (Figure~\ref{fig:capital}). Changing how
electricity is produced is therefore not the same as drawing investment into the sector.
That is the role of the investment incentives described in the next section.

\begin{figure}[H]
\centering
\includegraphics[width=\linewidth]{capital-intensity-sectors.pdf}
\caption{Generation capital intensity: long-run change in electricity output, capital and
labour input. Percent change from the control, Philippine example. Labour input is
effective hours, not a count of jobs.}
\label{fig:capital}
\end{figure}

\subsection{Air pollution and health}

Cleaner air reduces premature deaths and illness. OG--CLEWS takes the change in
fine-particle (PM$_{2.5}$) emissions from the resource plan and, using country data from
the Global Burden of Disease, adjusts survival rates and working capacity by age. OG-Core
then follows the resulting changes in population, labour supply, saving and consumption.

The two effects reach the economy in different ways. Most deaths from air pollution occur
late in life, near or after retirement, so fewer deaths add little to the labour force and
more to the number of pensioners. Illness also affects people of working age, and reduces
the time and effort they can give to work. Following both requires a model in
which the population has an age structure that changes over time. The economic effect can
be small even when the health benefit is large: GDP does not measure the value of longer
and healthier lives, and the two should be reported separately.

\pending{lives saved, and the effects on GDP and consumption, for the Philippine example
from the corrected health run.}
